\documentclass[12pt,letterpaper]{article}
\usepackage[utf8]{inputenc}
\usepackage[english]{babel}
\usepackage{biblatex}
\usepackage{amsmath,amsfonts,amssymb,amsthm}
\usepackage{mathtools}
\usepackage{listings, listings-rust}
\usepackage{xcolor}
\usepackage{mathpazo}
\usepackage[margin=2cm]{geometry}

\addbibresource{references.bib}

\newcommand*{\N}{\mathbb N}
\newcommand*{\code}[1]{\texttt {#1}}

\definecolor{codegreen}{rgb}{0,0.6,0}
\definecolor{codegray}{rgb}{0.5,0.5,0.5}
\definecolor{codepurple}{rgb}{0.58,0,0.82}
\definecolor{backcolour}{rgb}{0.95,0.95,0.92}

\lstdefinestyle{mystyle}{
    backgroundcolor=\color{backcolour},   
    commentstyle=\color{codegreen},
		% keywordstyle=\color{magenta},
    numberstyle=\tiny\color{codegray},
    stringstyle=\color{codepurple},
    basicstyle=\ttfamily\footnotesize,
    breakatwhitespace=false,         
    breaklines=true,                 
    captionpos=b,                    
    keepspaces=true,                 
    numbers=left,                    
    numbersep=5pt,                  
    showspaces=false,                
    showstringspaces=false,
    showtabs=false,                  
    tabsize=2
}

\lstset{language=Rust, style=mystyle}

\author{Alex Loomis}
\title{Notes}

\begin{document}

\maketitle

\section{Statement of Problem}

Consider a collection of units on a (connected component of a) grid,
along with a collection of (at least as many) targets.
We wish to find a collection of paths for the units to the targets that
\begin{itemize}
	\item minimizes the makespan (the maximum time to target),
	\item involves no unit conflicts (collisions).
\end{itemize}

In our setting, the units are rectangular, and all of the same size.
The units move at discrete time steps, indexed by $\N$.
Each time step, each unit may either wait at its current location,
or move one cell in any cardinal direction,
subject to the constraint to remain on the grid.

We will take the upper-leftmost cell that a unit occupies to be its origin.
Each grid cell has a (strictly) positive integer cost to enter,
and a unit whose origin enters a cell must wait that many time steps before moving again.

For example, in the following diagram\dots

\section{Overview of Approach}

To find these paths, we will begin by finding individual paths for each unit,
then will iteratively improve upon those paths.
This has the same structure as an A* search.

\begin{lstlisting}
	fn cbs(spec: Specification) -> HashMap<Path> {
		let mut open = Vec::new();
		let mut init = HashMap::new();
		for unit in spec.units {
			let path = find_path(unit, spec.targets);
			init.insert(unit, path);
		}
		open.push(make_node(init));
		loop {
			let node = choose_node(open);
			if node.conflicts.is_empty() {
				return node.paths;
			}
			for child in node.expand_node() {
				open.push(child);
			}
		}
	}
\end{lstlisting}

The biggest determinants of how many loops will be required are
what nodes we return from \code{choose\_node} and from \code{expand\_node}.
In general, we want these functions to return nodes that are close to being correct.
We will follow a conflict-based search (CBS) approach.
In CBS, we add constraints to eliminate conflicts,
until we have a conflict-free collection of paths.

\section{MDDs}

We will represent unit paths by multi-decision diagrams (MDDs).
These are \dots
We will use them to find sets of mutually exclusive states.

\section{Mutexes}

Our nodes consist of the following information:
\begin{itemize}
	\item a set of constraints,
	\item a minimal length MDD for each unit,
	\item the length of the longest MDD.
\end{itemize}

\appendix

\section*{References}

\nocite{*}
\printbibliography

\end{document}
